% =========================================================
% Gauges and Tagged Partitions — Notes
% Chapter: continuity
% Volume III · Analysis
% =========================================================

\subsection{Gauges and Tagged Partitions}
\label{sec:gauges-tagged-partitions}

\begin{tcolorbox}[title={Toolkit --- Gauges and Tagged Partitions}]
\begin{itemize}
  \item A \textbf{gauge} $\delta:I\to(0,\infty)$ is a variable-width
    tolerance function allowing the mesh to depend on position.
  \item A tagged partition $\dot{\mathcal{P}}$ is $\delta$-fine if
    each subinterval $I_i$ lies inside the $\delta(t_i)$-neighbourhood
    of its tag $t_i$: $\operatorname{Within}(x,t_i,\delta(t_i))$ for
    all $x\in[x_{i-1},x_i]$.
  \item \textbf{Cousin's theorem}: every gauge on $[a,b]$ admits a
    $\delta$-fine tagged partition.
  \item Gauge proofs of Boundedness, EVT, Location of Roots, and
    Heine--Cantor directly construct the required constants from a
    $\delta$-fine partition.
\end{itemize}
\end{tcolorbox}

% ---------------------------------------------------------
% Definition — Partition
% ---------------------------------------------------------

\begin{definition}[Partition]
\label{def:partition}
A \textbf{partition} of $I=[a,b]$ is $P=\{x_0,\ldots,x_n\}$ with
$a=x_0<x_1<\cdots<x_n=b$, dividing $I$ into $n$ subintervals
$[x_{i-1},x_i]$.
\end{definition}

\begin{remark*}[Standard quantified statement]
\[
  \operatorname{PartitionOfInterval}(P,[a,b])
  \;\iff\;
  \exists n\in\mathbb{N}:\;
  a=x_0<x_1<\cdots<x_n=b.
\]
\end{remark*}

% ---------------------------------------------------------
% Definition — Tagged Partition
% ---------------------------------------------------------

\begin{definition}[Tagged Partition]
\label{def:tagged-partition}
A \textbf{tagged partition} $\dot{\mathcal{P}}=\{([x_{i-1},x_i],t_i)\}$
is a partition $P$ together with tags $t_i\in[x_{i-1},x_i]$ for each $i$.
\end{definition}

\begin{remark*}[Standard quantified statement]
\[
  \operatorname{TaggedPartition}(\dot{\mathcal{P}},[a,b])
  \;\iff\;
  \operatorname{PartitionOfInterval}(P,[a,b])
  \;\wedge\;
  \forall i: t_i\in[x_{i-1},x_i].
\]
\end{remark*}

\begin{remark*}[Interpretation]
The tag $t_i$ is a distinguished sample point from the $i$-th
subinterval — the natural setting for Riemann sums via $f(t_i)
(x_i-x_{i-1})$.
\end{remark*}

% ---------------------------------------------------------
% Definition — Gauge
% ---------------------------------------------------------

\begin{tcolorbox}[title={Definition --- Gauge}]
\begin{definition}[Gauge]
\label{def:gauge}
A \textbf{gauge} on $I=[a,b]$ is $\delta:I\to(0,\infty)$.
\end{definition}
\end{tcolorbox}

\begin{remark*}[Standard quantified statement]
\[
  \delta\text{ a gauge on }I
  \;\iff\;
  \delta:I\to\mathbb{R} \;\wedge\; \forall t\in I:\delta(t)>0.
\]
\end{remark*}

\begin{remark*}[Interpretation]
A gauge generalizes constant mesh $\delta$. At each $t\in I$,
$\delta(t)$ specifies a local tolerance. The continuity of $f$ at
$t$ provides the local $\delta(t)$; the gauge stitches these local
conditions into a global partition.
\end{remark*}

% ---------------------------------------------------------
% Definition — delta-fine Partition
% ---------------------------------------------------------

\begin{definition}[$\delta$-Fine Partition]
\label{def:delta-fine-partition}
$\dot{\mathcal{P}}=\{([x_{i-1},x_i],t_i)\}$ is
\textbf{$\delta$-fine} if $\forall i:
[x_{i-1},x_i]\subseteq(t_i-\delta(t_i),t_i+\delta(t_i))$,
i.e., $\operatorname{Within}(x,t_i,\delta(t_i))$ for all
$x\in[x_{i-1},x_i]$.
\end{definition}

\begin{remark*}[Standard quantified statement]
\[
  \dot{\mathcal{P}}\;\delta\text{-fine}
  \;\iff\;
  \forall i,\;\forall x\in[x_{i-1},x_i]:\;
  \operatorname{Within}(x,t_i,\delta(t_i)).
\]
\end{remark*}

\begin{remark*}[Definition predicate reading]
\[
  \dot{\mathcal{P}}\;\delta\text{-fine}
  \;\colonequiv\;
  \forall i:\;
  x_{i-1}\geq t_i-\delta(t_i) \;\wedge\; x_i\leq t_i+\delta(t_i).
\]
\end{remark*}

\begin{remark*}[Interpretation]
Each subinterval must lie within the $\delta(t_i)$-ball around its
own tag. The continuity condition at $t_i$ then controls $f(x)$ in
terms of $f(t_i)$ for all $x\in[x_{i-1},x_i]$.
\end{remark*}

% ---------------------------------------------------------
% Lemma — Every Point Is Covered by a Tag
% ---------------------------------------------------------

\begin{lemma}[Every Point Is Covered by a Tag]
\label{lem:every-point-covered-by-tag}
If $\dot{\mathcal{P}}$ is $\delta$-fine on $[a,b]$ and $x\in[a,b]$,
then $\exists i:\operatorname{Within}(x,t_i,\delta(t_i))$.

\smallskip\noindent
\hyperref[prf:every-point-covered-by-tag]{\textit{Go to proof.}}
\end{lemma}

\begin{remark*}[Standard quantified statement]
\[
  \dot{\mathcal{P}}\;\delta\text{-fine},\;x\in[a,b]
  \;\Rightarrow\;
  \exists i:|x-t_i|\leq\delta(t_i).
\]
\end{remark*>

\begin{remark*}[Interpretation]
The bridge between a $\delta$-fine partition and a uniform bound.
For any $x\in I$, some subinterval contains $x$; that subinterval
lies in the $\delta(t_i)$-ball of its tag; so the continuity
condition at $t_i$ controls $f(x)$.
\end{remark*}

% ---------------------------------------------------------
% Theorem — Cousin's Theorem
% ---------------------------------------------------------

\begin{theorem}[Cousin's Theorem]
\label{thm:cousins-theorem}
Every gauge $\delta$ on $[a,b]$ admits a $\delta$-fine tagged
partition of $[a,b]$.

\smallskip\noindent
\hyperref[prf:cousins-theorem]{\textit{Go to proof.}}
\end{theorem}

\begin{remark*}[Standard quantified statement]
\[
  \forall\delta:[a,b]\to(0,\infty),\;
  \exists\operatorname{TaggedPartition}(\dot{\mathcal{P}},[a,b]):\;
  \dot{\mathcal{P}}\;\delta\text{-fine}.
\]
\end{remark*}

\begin{remark*}[Interpretation]
The gauge analogue of Heine--Borel: just as every open cover of
$[a,b]$ has a finite subcover, every gauge admits a $\delta$-fine
partition. Proof by bisection: if no $\delta$-fine partition existed,
bisect repeatedly; the intersection point $t$ has $\delta(t)>0$
whose neighbourhood contains some final interval —
contradiction.
\end{remark*}

% ---------------------------------------------------------
% Remark — Gauge Proofs of Classical Theorems
% ---------------------------------------------------------

\begin{remark*}[Gauge proofs of the four classical theorems]
Each proof follows the same pattern: define a gauge $\delta$ from
the continuity hypothesis, invoke Cousin's theorem for a $\delta$-fine
partition, apply the Covering Lemma to reduce to finitely many tags,
derive the global conclusion.

\medskip
\textbf{Boundedness.} At each $t$, continuity gives $\delta(t)$ with
$\operatorname{Within}(x,t,\delta(t))\Rightarrow|f(x)-f(t)|\leq 1$.
$\delta$-fine partition with tags $t_1,\ldots,t_n$; set $K=\max|f(t_i)|$.
Covering Lemma: $\forall x,\,\exists t_i:|x-t_i|\leq\delta(t_i)$, so
$|f(x)|\leq|f(x)-f(t_i)|+|f(t_i)|\leq 1+K$.

\medskip
\textbf{EVT.} Let $M=\sup f(I)$; assume $f(x)<M$ for all $x$. At
each $t$, continuity gives $\delta(t)$ with
$\operatorname{Within}(x,t,\delta(t))\Rightarrow f(x)<\tfrac{1}{2}(M+f(t))$.
$\delta$-fine partition produces $\widetilde{M}<M$ bounding $f$
everywhere — contradiction with definition of $M$.

\medskip
\textbf{Location of Roots.} Assume $f(t)\neq 0$ for all $t$. At
each $t$, continuity gives $\delta(t)$ preserving
$\operatorname{sgn}(f(t))$. A $\delta$-fine partition forces
consistent sign from $f(a)<0$ to $f(b)>0$ — contradiction.

\medskip
\textbf{Heine--Cantor.} Given $\varepsilon>0$: at each $t$, let
$\delta(t)$ satisfy $\operatorname{Within}(x,t,2\delta(t))
\Rightarrow|f(x)-f(t)|\leq\varepsilon/2$. $\delta$-fine partition
yields $\delta_\varepsilon=\min\{\delta(t_i)\}$. For
$|x-u|\leq\delta_\varepsilon$: find $t_i$ covering $x$;
$|u-t_i|\leq|u-x|+|x-t_i|\leq 2\delta(t_i)$; so
$|f(x)-f(u)|\leq\varepsilon$.
\end{remark*}
